% =====================================================================
%  LAB DECK -- companion to the theory deck for the
%  Column Generation & Branch-and-Price workshop
% =====================================================================
\documentclass[10pt,aspectratio=169]{beamer}

\usetheme{metropolis}
\usecolortheme{default}
\metroset{progressbar=frametitle, sectionpage=progressbar, subsectionpage=progressbar}
\setbeamertemplate{frame numbering}{%
  \insertshortauthor \;|\; \insertframenumber\,/\,\inserttotalframenumber}

\usepackage[utf8]{inputenc}
\usepackage[T1]{fontenc}
\usepackage{amsmath,amssymb,mathtools}
\usepackage{tikz}
\usepackage{tcolorbox}
\usepackage{booktabs}
\usepackage{listings}
\usepackage{xcolor}
\usepackage{array}
\usetikzlibrary{positioning,arrows.meta,shapes.geometric,calc,fit,backgrounds}
\tcbuselibrary{skins,breakable}

\definecolor{accent}{HTML}{C72A2A}
\definecolor{accent2}{HTML}{1F6FB2}
\definecolor{soft}{HTML}{F2F2F2}
\definecolor{codebg}{HTML}{F7F7F7}
\definecolor{codekw}{HTML}{0033B3}
\definecolor{codestr}{HTML}{067D17}
\definecolor{codecom}{HTML}{8C8C8C}

\lstdefinestyle{py}{
  language=Python,
  basicstyle=\ttfamily\scriptsize,
  keywordstyle=\color{codekw}\bfseries,
  stringstyle=\color{codestr},
  commentstyle=\color{codecom}\itshape,
  showstringspaces=false,
  breaklines=true,
  backgroundcolor=\color{codebg},
  frame=single,
  rulecolor=\color{black!20},
  numbers=none,
  xleftmargin=4pt,
  xrightmargin=4pt,
  framexleftmargin=2pt,
  framexrightmargin=2pt,
  aboveskip=2pt,
  belowskip=2pt,
}
\lstdefinestyle{sh}{
  basicstyle=\ttfamily\small,
  showstringspaces=false,
  breaklines=true,
  backgroundcolor=\color{codebg},
  frame=single,
  rulecolor=\color{black!20},
  xleftmargin=4pt,
  xrightmargin=4pt,
  aboveskip=2pt,
  belowskip=2pt,
}
\lstset{style=py}

\newtcolorbox{keybox}[1][Key idea]{
  colback=accent!8,colframe=accent,sharp corners,
  boxrule=0.5pt, left=4pt,right=4pt,top=3pt,bottom=3pt,
  fonttitle=\bfseries, title={#1}}
\newtcolorbox{warnbox}[1][Pitfall]{
  colback=yellow!10,colframe=orange!80!black,sharp corners,
  boxrule=0.5pt,left=4pt,right=4pt,top=3pt,bottom=3pt,
  fonttitle=\bfseries,title={#1}}
\newtcolorbox{exbox}[1][Exercise]{
  colback=green!5,colframe=green!50!black,sharp corners,
  boxrule=0.5pt,left=4pt,right=4pt,top=3pt,bottom=3pt,
  fonttitle=\bfseries\sffamily,title={#1}}
\newtcolorbox{tipbox}[1][Tip]{
  colback=accent2!8,colframe=accent2!70,sharp corners,
  boxrule=0.5pt,left=4pt,right=4pt,top=3pt,bottom=3pt,
  fonttitle=\bfseries,title={#1}}

\newcommand{\R}{\mathbb{R}}
\newcommand{\Z}{\mathbb{Z}}
\newcommand{\rmp}{\mathsf{RMP}}
\newcommand{\redcost}{\bar{c}}

\title{Column Generation \& Branch-and-Price}
\subtitle{Lab deck -- installation, running, exercises}
\author{Antoine Legrain}
\institute{Polytechnique Montr\'eal}
\date{May 20, 2026}

\begin{document}

\begin{frame}\titlepage\end{frame}

\begin{frame}{What this deck contains}
  \begin{itemize}\setlength{\itemsep}{4pt}
    \item Setup: Python virtual environment, dependencies, smoke check.
    \item File map of the \texttt{participants/} folder.
    \item How to run the smoke check on \texttt{toy.txt}.
    \item Exercises A through E, in execution order, with each \texttt{EX-X.k}
          hole listed explicitly so nothing is skipped.
    \item Troubleshooting recipes.
    \item Cheat sheet: HiGHS \& \texttt{rcspp} APIs.
  \end{itemize}
  \vfill
  Open this deck (on a second screen) \emph{while you code}. The theory is in
  \texttt{theory.pdf}.
\end{frame}

% =====================================================================
\section{1. Installation}
% =====================================================================

\begin{frame}[fragile]{What you need on your laptop}
  \begin{itemize}\small
    \item Python $\ge 3.10$ recommended (works on $\ge 3.8$).
    \item \texttt{pip} and \texttt{venv} (standard library on most systems).
    \item \alert{Windows users}: use WSL, Git Bash, or Anaconda --
          \texttt{venv} works the same way once Python is on \texttt{PATH}.
    \item Network access for the first install: HiGHS / NumPy / NetworkX /
          matplotlib come from PyPI, \texttt{rcspp} from TestPyPI.
    \item About 200~MB of disk for the venv.
  \end{itemize}
  \medskip
  \begin{tipbox}[macOS / Homebrew]
    If your system has multiple Pythons, prefer the Homebrew one
    (\texttt{brew install python}) over the one shipped by Apple.
  \end{tipbox}
\end{frame}

\begin{frame}[fragile]{Create the virtual environment}
  \textbf{Linux / macOS:}\hspace{1em}
  \texttt{\scriptsize cd participants/ \&\& python3 -m venv .venv \&\&
  source .venv/bin/activate}\\[2pt]
  \textbf{Windows PowerShell:}\hspace{1em}
  \texttt{\scriptsize cd participants/ ; python -m venv .venv ;
  .venv\char`\\Scripts\char`\\activate}

  \medskip
  Then, in both cases:
\begin{lstlisting}[style=sh,basicstyle=\ttfamily\scriptsize]
(.venv) $ python -m pip install --upgrade pip
(.venv) $ pip install -r requirements.txt
(.venv) $ pip install --index-url https://test.pypi.org/simple/ \
                      --extra-index-url https://pypi.org/simple \
                      rcspp==0.0.1.dev0
\end{lstlisting}
  \begin{warnbox}[Always activate the venv]
    Re-activate the venv in every new terminal, otherwise
    \texttt{import highspy} or \texttt{import rcspp} picks up the system
    Python.
  \end{warnbox}
\end{frame}

\begin{frame}[fragile]{Running against the reference solution}
  Once your environment is healthy you can run any of the same checks
  using the reference implementation under \texttt{solutions/}:
\begin{lstlisting}[style=sh,basicstyle=\ttfamily\scriptsize]
(.venv) $ cd solutions/
(.venv) $ python -m vrp.smoke_test
(.venv) $ python -m vrp.tests.exA_master
(.venv) $ python -m vrp.tests.exB_pricing
(.venv) $ python -m vrp.tests.exC_cg
(.venv) $ python -m vrp.tests.exD_diving
(.venv) $ python -m vrp.tests.exE_bnp
\end{lstlisting}
  Each check defaults to the small \texttt{toy.txt} instance; pass
  \texttt{--instance R101\_25.txt} for a larger Solomon file.
\end{frame}

\begin{frame}[fragile]{Smoke check -- you should see this}
\begin{lstlisting}[style=sh]
(.venv) $ python -m vrp.smoke_test
[smoke] highspy version: 1.10.0
[smoke] rcspp version:   0.0.1
[smoke] toy instance:    7 customers, capacity 50
[smoke] rcspp solve:     1 sink label, cost = -2.00
[smoke] OK -- environment is ready.
\end{lstlisting}
  If you see \emph{any} other line marked \texttt{[smoke] ERROR}, stop and
  fix the install before going further -- the exercises will not work
  otherwise.
  \medskip
  \begin{tipbox}[Reproducing the smoke check]
    The smoke check does not touch any code you'll modify, so you can
    re-run it at any time to confirm your environment is still healthy.
  \end{tipbox}
\end{frame}

% =====================================================================
\section{2. Project tour}
% =====================================================================

\begin{frame}[fragile]{The \texttt{participants/} folder}
\begin{lstlisting}[style=sh,basicstyle=\ttfamily\tiny]
participants/
+-- README.md, requirements.txt, theory.pdf, lab.pdf
+-- instances/                 toy.txt + Solomon files (+ duals/<inst>/iter_*.txt)
+-- solutions/                 reference solutions (do not peek!)
+-- vrp/
    +-- instance.py, instance_reader.py    (read-only)
    +-- smoke_test.py                      environment check
    +-- log.py                             vrp.log configure / get_logger
    +-- plot.py                            CG / B&P convergence figures
    +-- vrp.py                             HOLES: A.3, B.3, C.2-C.4
    +-- cg/
    |   +-- master_problem.py              HOLES: A.1, A.2, D.1
    |   +-- pricing.py                     HOLES: B.1, B.2, E.4
    |   +-- bound.py                       HOLE : C.5
    |   +-- diving.py                      HOLES: D.2, D.3
    |   +-- bnp.py                         HOLES: E.1, E.2, E.3 (loop is provided)
    |   +-- mp_solution.py, path.py        (read-only)
    +-- tests/                             progressive checks (NOT pytest)
        +-- exA_master.py, exB_pricing.py, exC_cg.py,
        +-- exD_diving.py, exE_bnp.py,
        +-- compare_cg.py, compare_bnp.py  parametric studies + matplotlib
\end{lstlisting}
\end{frame}

\begin{frame}{The TODO holes -- complete checklist}
  \tiny
  \begin{tabular}{@{}lll@{}}
    \toprule
    Tag & File & What to fill in \\
    \midrule
    \textbf{EX-A.1} & \texttt{cg/master\_problem.py} & vehicle constraint $\sum_p x_p\!\le\!K$, plus one $=\!1$ row per customer \\
    \textbf{EX-A.2} & \texttt{cg/master\_problem.py} & add one column per path (no upper bound!) \\
    \textbf{EX-A.3} & \texttt{vrp.py} & generate one round-trip per customer \\
    \midrule
    \textbf{EX-B.1} & \texttt{cg/pricing.py} & build the resource graph (nodes + arcs) \\
    \textbf{EX-B.2} & \texttt{cg/pricing.py} & per-arc base cost \& dual \texttt{Row} \\
    \textbf{EX-B.3} & \texttt{vrp.py} & wire \texttt{solve\_subproblem} (build once + \texttt{update\_reduced\_costs}) \\
    \midrule
    \textbf{EX-C.1} & \texttt{vrp.py} & multi-column add to master \\
    \textbf{EX-C.2} & \texttt{vrp.py} & Wentges dual smoothing (optional) \\
    \textbf{EX-C.3} & \texttt{vrp.py} & per-iteration log line (LP, LB, UB, gap) \\
    \textbf{EX-C.4} & \texttt{cg/bound.py} & Lagrangian bound formula \\
    \midrule
    \textbf{EX-D.1} & \texttt{cg/master\_problem.py} & implement \texttt{\_set\_integrality} (RMH) \\
    \textbf{EX-D.2} & \texttt{cg/diving.py} & pick the column to fix \\
    \textbf{EX-D.3} & \texttt{cg/diving.py} & re-run CG after a fix \\
    \midrule
    \textbf{EX-E.1} & \texttt{cg/bnp.py} & aggregated arc flow $\bar y_{uv}$ \\
    \textbf{EX-E.2} & \texttt{cg/bnp.py} & most-fractional arc \\
    \textbf{EX-E.3} & \texttt{cg/bnp.py} & build the two children \\
    \textbf{EX-E.4} & \texttt{cg/pricing.py} & skip forbidden arcs \\
    \bottomrule
  \end{tabular}
  \smallskip\\
  \scriptsize The B\&P search loop itself is \alert{already coded} in \texttt{bnp.py}.
\end{frame}

% =====================================================================
\section{3. Exercise A -- Master in HiGHS}
% =====================================================================

\begin{frame}[fragile]{Exercise A -- goal}
  \begin{exbox}[Exercise A -- master \& initial columns]
    \textbf{Holes to fill: EX-A.1, EX-A.2, EX-A.3.}
    \begin{itemize}\small
      \item \textbf{A.1}: in \texttt{cg/master\_problem.py}, build the
            \alert{vehicle-count constraint} $\sum_p x_p\!\le\!K$ \emph{and}
            one $=\!1$ covering row per customer. (The vehicle row gives
            the dual $\sigma$ that the Lagrangian bound needs.)
      \item \textbf{A.2}: in the same file, add a column per Path.
            Bounds: \texttt{lower=0}, \alert{no upper bound}
            (\texttt{upper = highspy.kHighsInf}). The column has
            coefficient~1 on the vehicle row.
      \item \textbf{A.3}: in \texttt{vrp.py}, complete
            \texttt{generate\_initial\_paths}: one round-trip
            $0\!\to\!i\!\to\!0$ per customer.
    \end{itemize}
  \end{exbox}
  \medskip
  \textbf{Run the check}
\begin{lstlisting}[style=sh]
(.venv) $ python -m vrp.tests.exA_master
\end{lstlisting}
\end{frame}

\begin{frame}[fragile]{Exercise A -- expected output}
\begin{lstlisting}[style=sh]
[exA_master] building master with 7 initial columns
[exA_master] LP optimum: 414.18  (expected ~ 414.18)
[exA_master] dual prices:
  pi_1=42.43  pi_2=41.23  pi_3=22.36
  pi_4=41.23  pi_5=50.00  pi_6=31.05
[exA_master] OK
\end{lstlisting}
  \begin{tipbox}[Why no upper bound on $x_p$?]
    The covering constraint $\sum_p a_{ip}\,x_p=1$ already bounds $x_p$
    from above whenever $a_{ip}\ge 1$. Adding $x_p\!\le\!1$ would create an
    extra dual $\mu_p$ that the pricer would have to subtract from
    $\redcost_p$. We never do that.
  \end{tipbox}
\end{frame}

% =====================================================================
\section{4. Exercise B -- Pricing wiring}
% =====================================================================

\begin{frame}[fragile]{Exercise B -- goal}
  \begin{exbox}[Exercise B -- arc costs \& RCSPP solve]
    \textbf{Holes to fill: EX-B.1, EX-B.2, EX-B.3.}
    \begin{itemize}\small
      \item \textbf{B.1}: in \texttt{cg/pricing.py}, build the graph
            \emph{once} -- 3 resources, all customer nodes, a sink, and
            all arcs.  No duals yet.
      \item \textbf{B.2}: per arc, push the \alert{base} resource
            increments $(d_{uv},\, s_u\!+\!d_{uv},\, q_v)$ and a dual
            \texttt{Row(origin\_id, 1.0)} so that the reduced cost
            becomes $\redcost_{uv}\!=\!d_{uv}-\pi_{u}$ after
            \texttt{update\_reduced\_costs}.
      \item \textbf{B.3}: in \texttt{vrp.py}, wire
            \texttt{solve\_subproblem(duals)}: lazily build the graph on
            the first call, then \texttt{update\_reduced\_costs(duals)}
            and return \texttt{graph.solve()} each iteration.
    \end{itemize}
  \end{exbox}
  \medskip
\begin{lstlisting}[style=sh]
(.venv) $ python -m vrp.tests.exB_pricing
\end{lstlisting}
\end{frame}

\begin{frame}[fragile]{Exercise B -- expected output}
\begin{lstlisting}[style=sh]
[exB_pricing] instance toy loaded: 6 customers, 6 vehicles, capacity 50
[exB_pricing] loading saved duals from instances/duals/R101/iter_0.txt
[exB_pricing] best route        : [0, 3, 1, 2, 7]
[exB_pricing] reduced cost      : -2999917.54
[exB_pricing] # neg-cost routes : 7
[exB_pricing] OK
\end{lstlisting}
\end{frame}

% =====================================================================
\section{5. Exercise C -- Efficient CG}
% =====================================================================

\begin{frame}{Exercise C -- goal}
  \begin{exbox}[Exercise C -- multi-column, smoothing, LB]
    \textbf{Holes to fill in \texttt{vrp.py},: EX-C.1, EX-C.2, EX-C.3, EX-C.4.}
    \begin{itemize}\small
      \item \textbf{C.2}: cap pricer output at \texttt{K\_MAX} columns; only
            those with $\redcost<-\varepsilon$ are added to the master.
      \item \textbf{C.3}: Wentges smoothing
            $\boldsymbol\pi\!\leftarrow\!\alpha\boldsymbol\pi^{prev}
                                     +(1-\alpha)\boldsymbol\pi$.
      \item \textbf{C.4}: per-iteration log line
            \texttt{iter, n\_cols, LP, LB, UB, gap}, plus the gap-stop
            (\texttt{gap\_pct}) and UB-cut (\texttt{incumbent\_ub}) early
            exits already wired in \texttt{solve\_cg}.
      \item \textbf{C.5}: in \texttt{cg/bound.py}, implement
            $\text{LB}\!=\!\sum_i\pi_i + K\sigma + K\!\cdot\!\min(0,\redcost^*)$
            (Trick~7 in the theory deck).
    \end{itemize}
    Try \texttt{--K 1\,5\,20\,50} and \texttt{-a 0\,0.3\,0.5}
    (or \texttt{compare\_cg} below) to feel each knob.
  \end{exbox}
  \begin{tipbox}[Don't rebuild the master or the graph]
    Build both \emph{once} in Exercise A/B. Each CG iteration calls
    \texttt{master.add\_column(\dots)} on the new columns and
    \texttt{graph.update\_reduced\_costs(\dots)} -- both warm-start.
  \end{tipbox}
\end{frame}

\begin{frame}[fragile]{Exercise C -- expected output (R101\_25)}
\begin{lstlisting}[style=sh,basicstyle=\ttfamily\scriptsize]
(.venv) $ python -m vrp.tests.exC_cg --instance R101_25.txt --K 50
iter cols  LP              LB              UB         gap%    min_rc    
0    75    10000602.89     -50997693.17    inf        119.610 -3999890.29
1    92    762.88          -346.80         inf        319.976 -73.98    
2    108   701.73          10.90           inf        6337.639 -46.06    
3    117   641.86          396.93          inf        61.708  -16.33    
4    122   641.86          248.90          inf        157.879 -26.20    
5    124   637.33          526.60          inf        21.028  -7.38     
6    125   637.33          536.92          inf        18.702  -6.69     
7    126   637.33          590.33          inf        7.961   -3.13     
8    127   636.28          636.24          inf        0.006   -0.00     
9    127   636.28          636.28          inf        0.000   -0.00   <-- converged
\end{lstlisting}
  Try $K\!\in\!\{1,5,20,100\}$ and observe the iteration count drop.
  Check what Wentges $\alpha\!=\!0.5$ does.
\end{frame}

% =====================================================================
\section{6. Exercise D -- Heuristic IP}
% =====================================================================

\begin{frame}[fragile]{Exercise D -- goal}
  \begin{exbox}[Exercise D -- RMH and plain diving]
    \textbf{Holes to fill: EX-D.1, EX-D.2, EX-D.3.}
    \begin{itemize}\small
      \item \textbf{D.1}: in \texttt{cg/master\_problem.py}, complete
            \texttt{\_set\_integrality(integer)}; \alert{this is required}
            for the RMH to actually solve a binary MIP. Without it,
            \texttt{master.solve(relax=False)} silently returns the LP
            solution.
      \item \textbf{D.2, D.3}: in \texttt{cg/diving.py}, pick the column
            with the largest fractional value, fix it to~1, then re-run CG.
            \texttt{plain\_dive} also takes an optional \texttt{ub}: when
            the LP at a dive node already exceeds \texttt{ub}, the dive
            aborts (saves work inside B\&P).
    \end{itemize}
  \end{exbox}
\begin{lstlisting}[style=sh]
(.venv) $ python -m vrp.tests.exD_diving --instance R101_25.txt
\end{lstlisting}
\end{frame}

\begin{frame}[fragile]{Exercise D -- expected output}
\begin{lstlisting}[style=sh]
[root]     LP=617.1
[RMH]      incumbent=632.4   time=0.6s
[dive 1]   fix p_42 (x=1.0)  CG -> LP=620.8
[dive 2]   fix p_55 (x=0.92) CG -> LP=623.1
...
[dive end] incumbent=622.1   time=12.4s
[summary]  RMH=632.4   dive=622.1   LP=617.1
\end{lstlisting}
  \begin{tipbox}[Why diving sometimes loses]
    On easy instances the root pool is already very rich and RMH wins.
    Diving wins when fixing forces the pricer to discover \emph{new}
    compatible columns the root never saw.
  \end{tipbox}
\end{frame}

% =====================================================================
\section{7. Exercise E -- Branch-and-Price}
% =====================================================================

\begin{frame}{Exercise E -- goal}
  \begin{exbox}[Exercise E -- arc-branching B\&P]
    \textbf{Holes to fill in \texttt{cg/bnp.py}: EX-E.1, E.2, E.3, E.4.}
    \begin{itemize}\small
      \item \textbf{E.1}: aggregated arc flow $\bar y_{uv}$.
      \item \textbf{E.2}: most-fractional arc.
      \item \textbf{E.3}: two children -- forbid / force. The
            \emph{force} branch uses \texttt{\_up\_compatible(p,u,v)}: it
            drops every path that leaves $u$ to some $w\!\ne\!v$ or
            arrives at $v$ from $w\!\ne\!u$.
      \item \textbf{E.4}: in \texttt{pricing.py}, skip forbidden arcs at
            graph-build time.
    \end{itemize}
    The B\&P search loop is \alert{already coded} (\texttt{cg/bnp.py})
    and returns a \texttt{BnPIncumbent} carrying \texttt{cost},
    \texttt{sol} and the \texttt{\{ub,lb,time\}\_history} traces used by
    \texttt{plot\_bnp\_progress}. The Lagrangian bound was implemented in
    Exercise~C (Trick~7). Dive at root, RMH every \texttt{--rmh-every}
    nodes; \texttt{--gap} stops within $\delta\%$;
    \texttt{--no-depth-first} switches to pure best-first.
  \end{exbox}
\end{frame}

\begin{frame}[fragile]{Exercise E -- expected output}
\begin{lstlisting}[style=sh]
(.venv) $ python -m vrp.tests.exE_bnp --instance R101_25.txt
[bnp]   open=1  inc=622.1  best_lb=617.1  gap=0.8%
[node 1]  branch on (4,12)  bar_y=0.61
[node 2]  forbid (4,12)  LP=621.2 LB=618.3 depth=1
[node 3]  force  (4,12)  LP=620.4 LB=620.4 (integer)  inc=620.4
[bnp]   open=1  inc=620.4  best_lb=618.3
[node 4]  forbid further; LB=620.4  prune
[bnp]   DONE in 5 nodes,  optimum=620.4
\end{lstlisting}
\end{frame}

% =====================================================================
\section{8. Comparison runners}
% =====================================================================

\begin{frame}[fragile]{Two scripts to compare configurations}
  After exercises~C and~E you have working CG and B\&P. The lab kit
  ships two helpers that re-run them with different knobs and plot the
  convergence side-by-side (\texttt{vrp.plot} uses \texttt{matplotlib}):

\begin{lstlisting}[style=sh]
# CG -- compare pricing batch sizes
python -m vrp.tests.compare_cg --instance RC101.txt --Ks 1 5 20 50
# CG -- compare Wentges smoothing
python -m vrp.tests.compare_cg --instance RC101.txt --alphas 0 0.2 0.5
# B&P -- compare gap thresholds
python -m vrp.tests.compare_bnp --instance R101.txt --gaps 0 0.5 2
# B&P -- depth-first vs best-first
python -m vrp.tests.compare_bnp --instance R101.txt --search both
\end{lstlisting}
  Add \texttt{--save figure.png} to write the figure to disk instead of
  displaying it.
\end{frame}

\begin{frame}{What the figures show}
  \begin{itemize}\small
    \item \texttt{plot\_cg\_convergence} -- LP \& LB per iteration, plus
          the duality gap, one curve per run. Reads
          \texttt{CGStats.lp\_history} and \texttt{CGStats.lb\_history}.
    \item \texttt{plot\_bnp\_progress} -- UB / LB and gap per node
          \emph{and} per wall-clock second, one curve per run. Reads
          \texttt{BnPIncumbent.\{ub,lb,time\}\_history}.
  \end{itemize}
  Both \texttt{CGStats} and \texttt{BnPIncumbent} record their history
  as the algorithm progresses -- no extra wiring needed.
  \begin{keybox}[Use the runners for the lab write-up]
    Pick one figure that surprises you (e.g.\ where best-first beats
    depth-first, or where Wentges $\alpha\!=\!0.5$ wins on iterations
    but loses on time) and add a paragraph of analysis to your hand-in.
  \end{keybox}
\end{frame}

% =====================================================================
\section{9. Troubleshooting}
% =====================================================================

\begin{frame}[fragile]{Five errors that always show up}
  \scriptsize
  \begin{tabular}{@{}p{0.30\textwidth}p{0.30\textwidth}p{0.36\textwidth}@{}}
    \toprule
    Symptom & Cause & Fix \\
    \midrule
    \texttt{ImportError: highspy} & venv not active &
       \texttt{source .venv/bin/activate} \\
    LP value $=0$ at iter 0 & forgot $b_i=1$ on covering rows &
       set RHS=1 in EX-A.1 \\
    pricer returns empty & $\pi_0$ not zeroed &
       force \texttt{pi[0]=0} before pricing \\
    LB $>$ LP early on & $\redcost^*$ from heuristic pricer &
       LB only valid with exact pricer \\
    B\&P loops forever & forgetting to filter parent pool &
       drop incompatible columns at child node \\
    \bottomrule
  \end{tabular}
  \medskip
  \begin{tipbox}[Get unstuck fast]
    Each saved \texttt{iter\_*.txt} in \texttt{instances/duals/R101/} is a
    one-iteration regression check. Diff your duals against them.
  \end{tipbox}
\end{frame}

% =====================================================================
\section{10. Cheat sheet}
% =====================================================================

\begin{frame}[fragile]{HiGHS in Python -- the calls you need}
\begin{lstlisting}
import highspy as hp
m = hp.Highs()
m.silent()                                    # no chatter

# row: equality constraint sum a_ip * x_p == 1
m.addRow(1.0, 1.0, 0, [], [])                 # empty row, fill later
# column: cost = c_p, lb=0, ub=+inf  (no upper bound!)
m.addCol(c_p, 0.0, hp.kHighsInf,
         len(idx), idx, val)                  # idx = row indices

m.run()                                       # solve LP
duals = list(m.getSolution().row_dual)        # row duals = pi
sol   = list(m.getSolution().col_value)       # primal x_p
m.changeColIntegrality(j, hp.HighsVarType.kInteger)  # for MIP / RMH
m.changeColBounds(j, 1.0, 1.0)                # fix column to 1 (B&P / dive)
\end{lstlisting}
\end{frame}

\begin{frame}[fragile]{\texttt{rcspp} cheat sheet}
\begin{lstlisting}[basicstyle=\ttfamily\tiny]
from rcspp.graph     import ResourceGraph
from rcspp._core.graph import Row
from rcspp.resource  import (
    AdditionExtensionFunction, ValueDominanceFunction,
    ValueCostFunction, TrivialFeasibilityFunction,
    MinMaxFeasibilityFunction,
    TimeWindowExtensionFunction, TimeWindowFeasibilityFunction,
)

g = ResourceGraph()
g.add_real_resource(ext, feas, cost, dom)         # once per resource
g.add_node(node_id, is_source, is_sink)
g.add_arc((dist, time, load), o, d, arc_id, dist, # base costs
          [Row(o, 1.0)])                          # subtract pi_o

g.update_reduced_costs(dual_by_id)                # cheap re-pricing
sols = g.solve()                                  # list[Solution]
# .cost, .path_node_ids, .path_arc_ids
\end{lstlisting}
  \alert{Build the graph once} (in \texttt{PricingGraph.build}); each CG
  iteration calls only \texttt{update\_reduced\_costs(dual\_by\_id)}.
\end{frame}

\begin{frame}{Time budget reminder}
  \begin{tabular}{@{}llr@{}}
    \toprule
    Exercise & Topic & Budget \\
    \midrule
    A & Master in HiGHS                & 15 min \\
    B & Pricing wiring (RCSPP)         & 25 min \\
    -- & \emph{break (15 min)}         &        \\
    C & Multi-col + smoothing          & 20 min \\
    D & RMH + plain diving             & 25 min \\
    E & Arc-branching B\&P + LB        & 35 min \\
    \bottomrule
  \end{tabular}
  \medskip
  Total $\approx$ 2~h hands-on, leaving 2~h for theory and questions.
\end{frame}

\begin{frame}{Done?}
  \begin{itemize}\small
    \item Bonus exercise \alert{F}: replace plain dominance by
          \alert{ng-routes} in \texttt{pricing.py} and measure the
          speed-up on \texttt{R101\_50.txt}.
    \item Submit your solutions branch -- a reference solution will be
          published at the end of the workshop.
  \end{itemize}
  \vfill
  \centering Have fun, and break the LP, not the dominance!
\end{frame}

\end{document}
